\section{Matrix-Tree 定理}

需要 \texttt{Z}（模意义整数，见杂项章节）与高斯消元求行列式 \texttt{Guass::det}。

\subsection{生成树计数}

无向图生成树个数 = 拉普拉斯矩阵（度数矩阵 $-$ 邻接矩阵）去掉任一行一列后的行列式。重边按权值累加。

\begin{lstlisting}
namespace MatrixTree {
    using MInt = Z;
    using std::vector;
    using std::array;
    MInt undirected(int n,vector<array<int,3>> E){
        vector K(n,vector(n,MInt(0)));
        for (auto [x,y,w] : E){
            K[x][y] -= w;
            K[y][x] -= w;
            K[x][x] += w;
            K[y][y] += w;
        }
        for (int i = 0;i < n;i++){
            K[i].pop_back();
        }
        K.pop_back();
        return Guass::det(K);
    }
    // 该版本为外向树版本，如果为内向树则需要边权取反
    MInt directed(int n,int rt,vector<array<int,3>> E){
        vector K(n,vector(n,MInt(0)));
        for (auto [x,y,w] : E){
            K[x][y] -= w;
            K[y][y] += w;
        }
        for (int i = 0;i < n;i++){
            K[i].erase(K[i].begin()+rt,K[i].begin()+rt+1);
        }
        K.erase(K.begin()+rt,K.begin()+rt+1);
        return Guass::det(K);
    }
};
\end{lstlisting}

\subsection{BEST 定理}

若有向强连通图中存在欧拉回路，那么它的欧拉回路个数为 $ec(G)=t^{root}(G,k)\cdot\prod_{i=1}^n(\deg(i)-1)!$，其中 $k$ 可以取 $1\sim n$ 的任意正整数。

\begin{lstlisting}
namespace MatrixTree {
    using std::vector;
    using std::array;
    MInt BEST(int n,vector<array<int,3>> E){
        if (E.empty()) return 1;
        vector<int> in(n),out(n);
        for (auto [x,y,w] : E){
            out[x] += w;
            in[y] += w;
        }
        for (int i = 0;i < n;i++) {
            if (in[i] != out[i]) return 0;
        }
        MInt prod = 1;
        int mx = 0;
        for (int i = 0;i < n;i++) mx = std::max(mx,in[i]);
        vector<MInt> fact(mx+1,1);
        for (int i = 1;i <= mx;i++) fact[i] = fact[i-1] * i;
        for (int i = 0;i < n;i++) if (in[i]) prod *= fact[in[i]-1];
        return directed(n,0,E) * prod;
    }
};
\end{lstlisting}

